% =====================================================================
%  Sezione 2.10 — Finestre temporali e divisione dei dati
%  Sostituisce la parte della vecchia sezione "Finestre temporali e
%  labeling" relativa alle finestre e allo split. L'etichettatura e'
%  trattata nella sezione successiva.
% =====================================================================

\section{Finestre temporali e divisione dei dati}
\label{sec:finestre}

Come discusso nella Sezione~\ref{sec:motion}, una caduta non è riconoscibile da
un singolo fotogramma: alcune attività quotidiane — sedersi rapidamente,
chinarsi a raccogliere un oggetto — producono in un istante una configurazione
simile a quella di una caduta, e la differenza sta nella traiettoria temporale.
Il classificatore lavora quindi su \textbf{finestre}, cioè su sequenze ordinate
di vettori di grandezze.

\subsection{Lunghezza della finestra}

Detta $f_t$ la rappresentazione al fotogramma $t$, una finestra è
$X_t = [\,f_t, f_{t+1}, \dots, f_{t+L-1}\,]$, con $L$ da scegliere.

Il compromesso è noto: una finestra corta osserva solo un frammento del
movimento e può non contenere abbastanza contesto per distinguere una caduta da
un'attività rapida; una finestra lunga contiene più contesto ma diluisce
l'evento fra molti fotogrammi normali, e aumenta la latenza decisionale perché
il sistema deve attendere più fotogrammi prima di classificare.

Il valore utilizzato è $L = 22$ campioni a 9 fotogrammi al secondo, pari a
$2{,}44$ secondi. La scelta merita due precisazioni.

La prima riguarda la sua origine. Il valore deriva da una scelta iniziale di 45
campioni a 18 fotogrammi al secondo, riportata alla frequenza di esercizio del
dispositivo mantenendo invariata la durata coperta. Quel 45 non era però
giustificato da alcuna misura, ed è stato successivamente sottoposto a
verifica.

La seconda riguarda l'esito di quella verifica, riportato nel
Capitolo~\ref{cap:sperimentale}. La durata mediana della fase di caduta è di
$3{,}11$ secondi, quindi superiore alla finestra: l'ipotesi ragionevole era che
allungarla migliorasse il risultato. Il confronto su quattro lunghezze — 22, 28,
36 e 44 campioni, da $2{,}4$ a $4{,}9$ secondi — mostra invece che
\textbf{nessuna differenza è statisticamente distinguibile}, e che l'andamento
non è nemmeno monotono. La spiegazione più probabile è che la memoria effettiva
del sistema non coincida con la finestra: ogni campione porta con sé le
grandezze integrali accumulate sui 31 fotogrammi precedenti
(Sezione~\ref{sec:integrali}), cioè $1{,}7$ secondi di storia aggiuntiva.

Il valore 22 non è quindi un'eredità: è un valore verificato contro tre
alternative.

\subsection{Passo e sovrapposizione}

Le finestre sono sovrapposte, con passo di 2 campioni. Il motivo è che una
caduta può iniziare in un punto qualsiasi della sequenza: con finestre disgiunte
l'evento potrebbe cadere a cavallo di due finestre e non essere contenuto per
intero in nessuna delle due.

La sovrapposizione introduce però una criticità di cui la sezione successiva si
occupa: due finestre consecutive condividono venti campioni su ventidue e sono
quindi quasi identiche.

\subsection{La divisione dei dati}

Questo è il punto in cui l'impostazione iniziale del lavoro conteneva un difetto
di protocollo, ed è opportuno descriverlo perché condiziona
l'interpretazione di tutti i risultati precedenti alla sua correzione.

Nella prima versione la divisione fra addestramento e verifica avveniva
\textbf{a livello di finestra}, con partizionamento casuale. Poiché le finestre
si sovrappongono, finestre quasi identiche finivano da entrambe le parti: il
modello veniva quindi valutato su sequenze praticamente già viste, e la stima
delle prestazioni era ottimistica per costruzione.

Una prima correzione ha spostato la divisione \textbf{a livello di video}, il
che elimina il problema della sovrapposizione ma non quello dell'identità del
soggetto. Con pochi soggetti, dividendo per video, la stessa persona compare
nell'insieme di addestramento e in quello di verifica: il modello non viene
valutato su persone nuove, ma su persone note in riprese diverse. Su un dataset
a cinque soggetti la misura è netta — tutti e cinque i soggetti risultavano
presenti da entrambe le parti.

La divisione adottata è quindi \textbf{per soggetto}, secondo lo schema
\emph{leave-one-subject-out}: in ciascuna ripetizione un soggetto costituisce
l'insieme di verifica, un secondo quello di validazione e i restanti quello di
addestramento, con i tre insiemi disgiunti per persona. La domanda che si
misura diventa quella che interessa: \emph{il sistema funziona su una persona
mai vista?}

Due precisazioni completano il protocollo.

\paragraph{Lo scaler.} I parametri di normalizzazione sono stimati sui soli
fotogrammi di addestramento della ripetizione corrente, mai sull'intero
dataset — altrimenti l'informazione sulla distribuzione dell'insieme di verifica
rientrerebbe dalla porta di servizio. Sono inoltre stimati sui
\emph{fotogrammi} e non sulle finestre: con passo 2 su finestre da 22, ogni
fotogramma comparirebbe in undici finestre e peserebbe quindi undici volte, con
un peso per giunta variabile agli estremi di ogni ripresa.

\paragraph{I buchi di posa.} Le finestre non attraversano mai il confine fra due
riprese, né un'interruzione della posa. Ventidue campioni presi a cavallo di
un'interruzione descriverebbero un movimento mai avvenuto, e la finestra
coprirebbe un intervallo di tempo maggiore di quello che dichiara. Ogni ripresa
viene quindi suddivisa nei suoi tratti continui di posa valida, e le finestre si
costruiscono all'interno di un singolo tratto.

\subsection{La soglia decisionale}

Un'ultima componente del protocollo riguarda la soglia oltre la quale la
probabilità prodotta dal modello viene interpretata come allarme.

Nella versione iniziale la soglia veniva scelta sullo stesso insieme su cui poi
si misurava il risultato, il che equivale ad adattarsi a quei dati: il numero
che ne esce non è una previsione di come il sistema si comporterà altrove. La
soglia va quindi scelta su un terzo insieme — il validation — e applicata alla
verifica una volta sola.

Anche questa correzione, però, si è rivelata insufficiente, e la ragione è
riportata nel Capitolo~\ref{cap:sperimentale}. Con un solo soggetto di
validazione la stima della soglia è talmente rumorosa da essere
controproducente: le soglie selezionate nelle dieci ripetizioni coprivano
l'intero intervallo utile, da $0{,}01$ a $0{,}99$, con deviazione standard
$0{,}37$. Il confronto fra sei regole di selezione mostra che una \textbf{soglia
fissa, dichiarata a priori}, ottiene sia il risultato medio migliore sia la
dispersione minore, semplicemente perché non stima nulla dai dati e quindi non
può sovradattarsi.

La soglia utilizzata è quindi $0{,}50$, fissata prima di osservare i dati e
identica per tutte le ripetizioni e tutte le configurazioni. Il validation
conserva il proprio ruolo nell'arresto dell'addestramento, non nella decisione.
